\documentclass[10pt,a4paper]{article}
\usepackage[T1]{fontenc}
\usepackage[utf8]{inputenc}
\usepackage[portuguese]{babel}
\usepackage[a4paper,margin=1.25cm]{geometry}
\usepackage[hidelinks]{hyperref}
\usepackage{enumitem}
\usepackage{titlesec}
\usepackage{xcolor}
\definecolor{graytext}{gray}{0.35}
\definecolor{accent}{HTML}{1A4F8B}
\pagestyle{empty}
\setlength{\parindent}{0pt}
\setlist[itemize]{leftmargin=1em,itemsep=0.5pt,topsep=1pt,parsep=0pt}
\titleformat{\section}{\normalsize\bfseries\color{accent}}{}{0pt}{}[\vspace{-3pt}\rule{\linewidth}{0.4pt}]
\titlespacing*{\section}{0pt}{6pt}{6pt}
\newcommand{\cventry}[3]{\textbf{#1}\ \textcolor{graytext}{· #2}\hfill\textcolor{graytext}{\small #3}\par}
\newcommand{\tagline}[1]{\textcolor{graytext}{\small #1}\par}
\linespread{0.96}\selectfont
\begin{document}

{\LARGE\bfseries Mohamed Rafaa TOUNEKTI}\hfill{\small DevOps \& DevSecOps}\par
\vspace{3pt}
{\footnotesize \href{https://linkedin.com/in/mohamed-rafaa-tounekti}{linkedin.com/in/mohamed-rafaa-tounekti} \textcolor{graytext}{·} \href{https://wa.me/21622304599}{+216 22 304 599} \textcolor{graytext}{·} \href{tel:+5511929017154}{+55 11 92901-7154} \textcolor{graytext}{·} \href{mailto:rafaabou@gmail.com}{rafaabou@gmail.com} \textcolor{graytext}{·} Baseado em São Paulo (CPF Pronto) – Brasil \textcolor{graytext}{·} \href{https://github.com/rafaabou}{github.com/rafaabou} \textcolor{graytext}{·} \href{https://devops-devsecops.kdns.fr/}{devops-devsecops.kdns.fr}}\par
\vspace{2pt}\hrule\vspace{2pt}
\section*{Perfil}
Engenheiro DevOps e DevSecOps focado na construção de pipelines CI/CD confiáveis, na segurança de clusters Kubernetes e na automação de infraestrutura em escala. Atuo com AWS, sistemas Linux e ambientes conteinerizados, com experiência aprofundada em GitLab CI/CD, Terraform, Ansible e stacks de monitoramento.\par
\section*{Competências Técnicas}
\textbf{Cloud}: AWS (EC2, EKS, ECR, IAM, S3, VPC, Lambda)\par
\textbf{Contêineres \& Orquestração}: Docker, Kubernetes, Helm, Kustomize, Istio, OPA, Kyverno\par
\textbf{CI/CD \& GitOps}: GitLab CI/CD, Jenkins, ArgoCD, Flux CD, GitLab Runners\par
\textbf{Infraestrutura como Código}: Terraform, Ansible\par
\textbf{Segurança}: Trivy, GitLeaks, Semgrep, OWASP ZAP, SonarQube, Checkov, KubeLinter, Falco, HashiCorp Vault, AWS Secrets Manager\par
\textbf{Monitoramento \& Logging}: Prometheus, Grafana, Loki, Alertmanager\par
\textbf{Redes · Scripts · Banco de Dados · Armazenamento}: Nginx Ingress, Traefik, Cert-Manager, External-DNS, Bash, Python, PostgreSQL, MongoDB, Redis, NFS, Harbor, Nexus\par
\textbf{Virtualização · Ferramentas de IA}: KVM, VirtualBox, Vagrant, MCP, OpenCode, AI-assisted infra\par
\section*{Experiência}
\cventry{Engenheiro DevOps \& DevSecOps}{TUNISIE SOFTWARE FACTORY \& BE-SOFTILYS}{Jun 2025 – Out 2024}
\begin{itemize}
  \item Construí pipelines GitLab CI/CD com escaneamento Trivy, GitLeaks e Semgrep. Imagens vulneráveis bloqueadas antes do staging
  \item Implantei segurança zero-trust nos clusters Kubernetes com Istio mTLS, NetworkPolicies e OPA Gatekeeper
  \item Integrei o HashiCorp Vault aos fluxos de implantação com rotação automática de segredos via External Secrets Operator
  \item Configurei dashboards Prometheus \& Grafana com roteamento pelo Alertmanager; MTTR caiu significativamente
  \item Configurei HPA e PodDisruptionBudgets para suportar picos de tráfego sem perder requisições
  \item Automatizei o provisionamento de certificados TLS com Cert-Manager para todos os recursos de ingress
  \item Executei testes de validação de infraestrutura de ponta a ponta após cada ciclo de provisionamento; detectei desvios antes da produção
  \item Apliquei gates de segurança no CI/CD com limites de vulnerabilidade; pipelines falhavam quando os limites eram excedidos
  \item Centralizei o rastreamento de vulnerabilidades no DefectDojo para acompanhar a remediação
  \item Depurei problemas de runtime no Kubernetes — CrashLoopBackOff, ImagePullBackOff, falhas de agendamento
  \item Escrevi runbooks e documentação de implantação; relatórios de validação mantiveram o histórico
  \item Participei de discussões pré-venda com clientes potenciais sobre opções de arquitetura
  \item Configurei servidores MCP para assistentes de IA consultarem dados de Kubernetes, CI/CD e monitoramento
  \item Incorporei assistentes de IA aos fluxos diários para análise de logs, triagem de incidentes e diagnóstico
  \item Usei ferramentas de IA para acelerar o desenvolvimento em Terraform e Ansible; menos erros de configuração
  \item Criei regras de proteção de branch com revisões obrigatórias e verificações automatizadas de pipeline
  \item Conectei o Alertmanager ao Jira; alertas geram tickets de serviço automaticamente
  \item Segurança da cadeia de suprimentos: geração de SBOM, assinatura de imagens com Cosign e rastreamento de dependências
  \item Implantei o Falco para segurança em runtime nas cargas de trabalho Kubernetes
\end{itemize}
\tagline{Kubernetes · Docker · Helm · ArgoCD · GitLab CI/CD · Terraform · AWS · GitLeaks · Semgrep · Trivy · OPA Gatekeeper · Kyverno · DefectDojo · Prometheus · Grafana · Istio · HashiCorp Vault · Harbor}
\vspace{1pt}
\vspace{3pt}
\cventry{Engenheiro DevOps}{DOPIFY SOFTWARE \& TECHNOLOGIES}{Abr 2024 – Jan 2024}
\begin{itemize}
  \item Projetei pipelines CI/CD para várias equipes com SAST, SCA e escaneamento de segredos em cada build
  \item Automatizei o provisionamento de servidores em dev, staging e produção com Ansible; eliminei desvios manuais
  \item Configurei o Flux CD para infraestrutura orientada a GitOps; ambientes declarativos e auditáveis
  \item Implantei o Grafana Loki para agregação centralizada de logs; análise de causa raiz mais rápida
  \item Implantações sem downtime com estratégias blue-green e canary usando overlays Kustomize
  \item Configurei o Nginx Ingress Controller com terminação TLS para roteamento externo
  \item Provisionei runners GitLab auto-hospedados com executor Docker; filas de CI diminuíram
  \item O Alertmanager encaminhou incidentes para Slack e e-mail; as pessoas certas acionadas no momento certo
  \item Automatizei scripts de teste para validar o estado da infraestrutura após o provisionamento
  \item Gerenciei a governança dos pipelines e o planejamento de cutover para releases de produção
  \item Rastreei degradação de desempenho por meio de logs de sistema e aplicação
  \item Documentei a arquitetura de infraestrutura, procedimentos de implantação e problemas conhecidos
  \item Proteção de branches Git com revisões obrigatórias e status checks; gates de qualidade antes do merge
\end{itemize}
\tagline{AWS · Kubernetes · Docker · Terraform · Ansible · Flux CD · Grafana · Prometheus · GitLab CI/CD · HashiCorp Vault · Trivy · SonarQube · Kustomize · Nginx · Alertmanager}
\vspace{1pt}
\vspace{3pt}
\cventry{Engenheiro DevOps}{GEO-TRACKING GROUP}{Out 2023 – Abr 2022}
\begin{itemize}
  \item Reconstruí pipelines CI/CD; tempo de implantação de dias para horas em mais de 15 microsserviços
  \item Reduzi o tamanho das imagens Docker em 60\% com builds multi-stage; tempos de pull e custos de armazenamento caíram
  \item Provisionei infraestrutura AWS com módulos Terraform; versionada e reproduzível
  \item Integrei o SonarQube para verificações automatizadas de qualidade de código e acompanhamento de débito técnico
  \item Escrevi regras de alerta no Prometheus para detecção proativa; quedas não planejadas caíram 40\%
  \item Automatizei atualizações de dependências e patches de segurança via GitOps em mais de 200 dependências
  \item Gerenciei imagens Docker e manifests Kubernetes no Nexus como fonte única de verdade
  \item Automatizei rotinas de backup de banco de dados com recuperação point-in-time para PostgreSQL e MySQL
  \item Reforcei a segurança dos pods com SecurityContext e segmentação de rede via IPTables/UFW
  \item Criei bibliotecas compartilhadas do Jenkins reutilizáveis para padronizar templates de pipeline
  \item Executei testes de instalação e upgrade; documentei resultados e sinalizei desvios
  \item Monitorei capacidade de armazenamento e métricas de desempenho; ajustei configurações para disponibilidade
  \item Respondi a incidentes e escrevi relatórios pós-incidente com etapas de remediação
  \item Integrei o LDAP para autenticação centralizada
\end{itemize}
\tagline{Jenkins · Docker · Kubernetes · Terraform · Ansible · ArgoCD · SonarQube · Nexus · Prometheus · Grafana · Python · Bash · AWS · PostgreSQL · Cert Manager}
\vspace{1pt}
\vspace{3pt}
\cventry{Administrador de Sistemas}{DIVERSAS ORGANIZAÇÕES}{Jan 2022 – Fev 2016}
\begin{itemize}
  \item Administrei servidores Linux (RHEL, Ubuntu, Debian) — patches, atualizações de segurança e tuning
  \item Automatizei o provisionamento de VMs com Ansible; tempo de configuração de novos ambientes diminuiu
  \item Administrei roteadores, switches e firewalls; mantive políticas de rede
  \item Configurei compartilhamento de arquivos NFS para acesso centralizado ao armazenamento entre equipes
  \item Configurei DNS, DHCP e serviços de rede internos
  \item Projetei e testei procedimentos de recuperação de desastres com validação regular de backups
  \item Administrei MySQL e PostgreSQL — backup, recuperação e otimização de desempenho
  \item Configurei regras de firewall IPTables e UFW para segmentar a rede e restringir acessos
  \item Implantei servidores web Apache e Nginx para aplicações internas e externas
  \item Integrei o Active Directory para autenticação centralizada em Linux/Windows
  \item Gerenciei contas de usuários, chaves SSH e políticas de sudo em ambientes corporativos
  \item Prestei suporte L2/L3 com análise de causa raiz
  \item Auxiliei usuários finais com problemas de hardware, software e rede; conduzi treinamentos
  \item Mantive a documentação de processos de TI, configurações de rede e troubleshooting
  \item Configurei sistemas de armazenamento sobre IPv4 com endereçamento e roteamento adequados
  \item Trabalhei com equipes técnicas para planejar melhorias de infraestrutura e automação
  \item Automatizei fluxos de alerta para ticket com resolução gerenciada por SLA
\end{itemize}
\tagline{Linux · DHCP · DNS · Active Directory · MySQL · PostgreSQL · VMware · KVM · Ansible · Bash · VPN · Firewalls · NFS · Apache · Nginx}
\vspace{1pt}
\section*{Formação}
\textbf{Bacharelado}\par
\tagline{Tecnologia da Informação (Redes e Serviços de Computação) · I.S.E.T. · 2011 – 2014}
\textbf{Ensino Médio}\par
\tagline{Ciência da Computação · 2010 – 2011}
\section*{Idiomas}
Árabe (Nativo), Francês (Profissional), Inglês (Profissional)\par
\section*{Principais Conquistas}
\begin{itemize}
  \item Arquitetura de segurança zero-trust no Kubernetes
  \item Escaneamento automatizado de segurança nos pipelines CI/CD
  \item GitOps com ArgoCD \& Flux CD
  \item Alta disponibilidade com PDB e escalonamento automático
  \item Stack completa de monitoramento e alertas
\end{itemize}
\vspace{6pt}\hrule\vspace{3pt}
{\footnotesize\color{graytext} Análise de logs assistida por IA · triagem de incidentes · diagnóstico de infraestrutura · servidores MCP\par}
\end{document}